\documentclass{article}
\usepackage[UTF8]{ctex}
\usepackage{amsmath,amssymb,amsfonts}
\usepackage{enumitem}
\usepackage{geometry}
\geometry{a4paper, margin=1in}
% === 错题本排版优化（不参与题目标识） ===
\setlength{\parindent}{0pt}
\setlength{\parskip}{0.35em plus 0.10em minus 0.08em}
\setlength{\emergencystretch}{3em}
\setlist{itemsep=0.25em, topsep=0.40em, parsep=0pt}
\allowdisplaybreaks[2]
% === 错题本排版优化结束 ===
\begin{document}
\section*{中国矿业大学2025年高等代数真题}
\section*{20．中国矿业大学2025年高等代数试卷}
一．填空题，每题5分，共50分。\\

\subsection{题目1}
1．多项式 \(4 x^{4}-7 x^{2}-5 x-1\) 的有理根是 \(\_\_\_\_\)．

\subsection{题目2}
2．\(f(x)=\left|\begin{array}{cccc}x & 1 & 2 & 3 \\ 1 & x & 1 & -1 \\ 2 & 1 & x & 1 \\ 3 & -1 & 2 x & 1\end{array}\right|\) ，则 \(f(x)\) 中 \(x^{3}\) 的系数是 \(\_\_\_\_\)．

\subsection{题目3}
3．设向量组 \(\alpha_{1}, \alpha_{2}, \cdots, \alpha_{n}\) 线性无关，判断向量组 \(\alpha_{1}+\alpha_{2}, \alpha_{2}+\alpha_{3}, \cdots, \alpha_{n-1}+\alpha_{n}, \alpha_{n}+\alpha_{1}\) 的线性相关性 \(\_\_\_\_\)．

\subsection{题目4}
4．\(A=\left(\begin{array}{cc}\frac{3}{2} & -\frac{1}{2} \\ \frac{1}{2} & \frac{1}{2}\end{array}\right)\) ，则 \(A^{2025}=\)

\subsection{题目5}
5．设 \(A\) 为 \(n\) 阶实可逆矩阵，且满足 \(A\) 与 \(-A\) 合同，则 \(\_\_\_\_\)．

\begin{itemize}
  \item[A．] \(n\) 一定是奇数．
  \item[B．] \(n\) 一定是偶数。
  \item[C．] \(n\) 可能是奇数也可能是偶数。
  \item[D．] \(\frac{n-1}{2}\) 和 \(\frac{n}{2}\) 中有一个为偶数。
\end{itemize}

\subsection{题目6}
6．设 \(A\) 是8阶半正定矩阵，\(A\) 的秩为3，则实向量空间 \(\mathbb{R}^{8}\) 的子空间 \(W=\left\{X \in \mathbb{R}^{8} \mid X^{\mathrm{T}} A X=0\right\}\) 的维数等于 \(\_\_\_\_\)。

\subsection{题目7}
7．设 \(A\) 是 4 阶方阵 \(|2 E+A|=0, A A^{\mathrm{T}}=2 E,|A|>0\) ，求伴随矩阵 \(A^{*}\) 的两个特征值 \(\_\_\_\_\)\\
\(\_\_\_\_\)．．\\
\(\_\_\_\_\)

\subsection{题目8}
8．设 \(\sigma\) 是实向量空间 \(\mathbb{R}^{n}\) 上的线性变换，满足 \(\sigma\left(x_{1}, x_{2}, \cdots, x_{n}\right)=\left(0, x_{1}, \cdots, x_{n-1}\right)\) ，则 \(\sigma\) 的核的维数及 \(\sigma\) 的一组基分别是 \(\_\_\_\_\)。

\subsection{题目9}
9．设 \(A, B\) 都为 \(n\) 阶正交矩阵，且 \(|A| \neq|B|\) ，则 \(|A+B|=\) \(\_\_\_\_\)．

\subsection{题目10}
10．\(A=\left(\begin{array}{ccc}1 & 3 & 2 \\ -2 & -6 & -4 \\ 3 & 9 & 6\end{array}\right), A\) 的最小多项式是 \(\_\_\_\_\)

二．（10 分）设 \(f(x)\) 是一个整系数多项式，证明：若存在一个偶数 \(a\) 和奇数 \(b\) 使得 \(f(a)\) 和 \(f(b)\) 都是奇数，则 \(f(x)\) 没有整数解．

三．（10 分）已知 \(A\) 为 \(n\) 阶矩阵，\(n \geq 2, A^{*}\) 是 \(A\) 的伴随矩阵，若存在 \(n\) 维非零列向量 \(\alpha\) 满足 \(A \alpha=0\)，证明：非齐次线性方程组 \(A^{*} X=\alpha\) 有解的充分必要条件是 \(A\) 的秩为 \(n-1\) 。

四．（15 分）已知 \(A\) 为 \(n\) 阶可逆矩阵，\(B\) 为 \(n\) 阶幂零矩阵（即存在正整数 \(k\) ，使得 \(B^{k}=O\) ），若\(A B=B A\) ，证明：\(A+B\) 可逆，并求 \(A+B\) 的逆．

五．（15 分）已知 \(B\) 为 \(n\) 阶正定矩阵，\(C\) 是 \(n \times m\) 列满秩矩阵，记 \(A=\left(\begin{array}{cc}B & C \\ C^{\mathrm{T}} & O\end{array}\right)\) ，证明：二次型\(f(X)=X^{\mathrm{T}} A X\) 的正，负惯性指数分别为 \(n, m\) 。

六．（15 分）设 \(A\) 是数域 \(P\) 上的 \(n\) 阶方阵，\(f(x), g(x)\) 是数域 \(P\) 上的两个多项式，且有 \((f(x), g(x))=\) 1 ，记 \(V, V_{1}, V_{2}\) 分别是齐次线性方程组 \(B C X=0, B X=0, C X=0\) 的解空间，其中 \(B=f(A), C=\) \(g(A)\) ，证明：\(V=V_{1} \oplus V_{2}\) ．

七．（15 分）已知 \(A\) 是 3 阶正交矩阵，且 \(|A|=1\) ．

\begin{itemize}
  \item[（1）] 证明： 1 是 \(A\) 的一个特征值．
  \item[（2）] 证明：\(A\) 的特征多项式形如 \(f(\lambda)=\lambda^{3}-a \lambda^{2}+a \lambda-1\) ，其中 \(a\) 是一个常数．
  \item[（3）] 若 \(A\) 的特征值都是实数，且 \(|A+E| \neq 0\) ，证明：\(A^{\mathrm{T}}=A^{2}-3 A+3 E_{3}\) ，这里 \(E_{3}\) 表示 3 阶单位矩阵。
\end{itemize}

八．（20 分）已知 \(\sigma\) 是 \(n\) 维线性空间 \(V\) 上的线性变换．\\
（1）已知 \(\lambda\) 是 \(\sigma\) 的特征值，叙述 \(\lambda\) 的几何重数的定义。\\
（2）证明：存在 \(\alpha \in V\) ，使得向量组 \(\alpha, \sigma(\alpha), \cdots, \sigma^{n-1}(\alpha)\) 是 \(V\) 的一组基当且仅当 \(\sigma\) 的任意特征值 \(\lambda\)的几何重数都等于1。
\end{document}
